% ============================================================================
% LANGENTWURF TIER 1: parecer y conocer
% ============================================================================
% Klasse 9 | Qué pasa 2, Unidad 1 | S. 18-21
% Fachfarbe: #c41e3a (Karminrot)
% ============================================================================

\documentclass[11pt,a4paper]{article}

% ============================================================================
% PACKAGES
% ============================================================================
\usepackage[utf8]{inputenc}
\usepackage[T1]{fontenc}
\usepackage[ngerman]{babel}
\usepackage{geometry}
\usepackage{fancyhdr}
\usepackage{lastpage}
\usepackage{xcolor}
\usepackage{tcolorbox}
\usepackage{enumitem}
\usepackage{tabularx}
\usepackage{longtable}
\usepackage{array}
\usepackage{booktabs}
\usepackage{hyperref}
\usepackage{amssymb}

% ============================================================================
% KONFIGURATION
% ============================================================================
\newcommand{\plantier}{1}
\newcommand{\fachid}{spanisch}
\newcommand{\fach}{Spanisch}
\newcommand{\fachfarbe}{c41e3a}

% --- Metadaten ---
\newcommand{\jahrgangsstufe}{9}
\newcommand{\schulart}{Gesamtschule}
\newcommand{\stundenthema}{parecer y conocer}
\newcommand{\reihenthema}{Unidad 1: La nueva familia}
\newcommand{\stundennummer}{01}
\newcommand{\reihenstunden}{4}
\newcommand{\gerniveau}{A2}
\newcommand{\lehrwerk}{Qué pasa 2}
\newcommand{\lehrwerkseite}{S. 18--21}

% --- Lehrkraft ---
\newcommand{\lehrkraft}{N.N.}
\newcommand{\ausbildungsstand}{Lehrkraft}

% --- Lerngruppe ---
\newcommand{\lerngruppengroesse}{24}
\newcommand{\maennlich}{12}
\newcommand{\weiblich}{12}
\newcommand{\divers}{0}

% ============================================================================
% LAYOUT
% ============================================================================
\geometry{left=2.5cm, right=2.5cm, top=2.5cm, bottom=2.5cm}
\definecolor{fach}{HTML}{\fachfarbe}
\definecolor{gray}{HTML}{666666}
\definecolor{lightgray}{HTML}{f5f5f5}
\definecolor{successgreen}{HTML}{2a7a4b}
\definecolor{warningorange}{HTML}{b35c00}
\definecolor{basisgreen}{HTML}{2a7a4b}
\definecolor{standardblue}{HTML}{2c5aa0}
\definecolor{erweiterungpurple}{HTML}{7b1fa2}

\tcbuselibrary{skins,breakable}

\newtcolorbox{infobox}[1][]{
    colback=lightgray,
    colframe=fach,
    fonttitle=\bfseries,
    title=#1,
    breakable
}

\newtcolorbox{scorebox}{
    colback=white,
    colframe=fach,
    fonttitle=\bfseries,
    title=Didaktik-Score,
    breakable
}

% Kopf-/Fußzeile
\pagestyle{fancy}
\fancyhf{}
\fancyhead[L]{\small\textcolor{gray}{\fach{} -- Jg. \jahrgangsstufe{} -- \stundenthema}}
\fancyhead[R]{\small\textcolor{gray}{Langentwurf Tier 1}}
\fancyfoot[C]{\small\thepage{} / \pageref{LastPage}}
\renewcommand{\headrulewidth}{0.4pt}
\renewcommand{\footrulewidth}{0pt}

% ============================================================================
% DOKUMENT
% ============================================================================
\begin{document}

% ============================================================================
% DECKBLATT
% ============================================================================
\begin{titlepage}
    \centering
    \vspace*{2cm}

    {\large\textcolor{gray}{\schulart}}\\[2cm]

    {\Huge\textcolor{fach}{\textbf{Langentwurf}}}\\[0.5cm]
    {\large Tier 1 -- Vollständig}\\[2cm]

    {\LARGE\textbf{\stundenthema}}\\[0.5cm]
    {\large im Rahmen der Unterrichtsreihe}\\[0.3cm]
    {\Large\textit{\reihenthema}}\\[2cm]

    \begin{tabular}{rl}
        \textbf{Fach:} & \fach{} (\gerniveau) \\[0.3cm]
        \textbf{Klasse:} & \jahrgangsstufe \\[0.3cm]
        \textbf{Lehrwerk:} & \lehrwerk, \lehrwerkseite \\[0.3cm]
        \textbf{Datum:} & \underline{\hspace{3cm}} \\[0.3cm]
        \textbf{Zeit:} & \underline{\hspace{3cm}} (Doppelstunde, 90 Min.) \\[0.3cm]
        \textbf{Raum:} & \underline{\hspace{3cm}} \\[0.3cm]
    \end{tabular}

    \vfill

    {\large Lehrkraft: \lehrkraft}\\[0.3cm]
    {\normalsize\textcolor{gray}{\ausbildungsstand}}\\[1cm]

\end{titlepage}

% ============================================================================
% INHALTSVERZEICHNIS
% ============================================================================
\tableofcontents
\newpage

% ============================================================================
% 1. BEDINGUNGSANALYSE
% ============================================================================
\section{Bedingungsanalyse}

\subsection{Lerngruppe}
\begin{tabular}{ll}
    Kursgröße: & \lerngruppengroesse{} SuS (\maennlich{} m / \weiblich{} w / \divers{} d) \\
    Jahrgangsstufe: & \jahrgangsstufe \\
    Sprachniveau: & \gerniveau{} (GER) \\
    Fremdsprache: & 2. Fremdsprache (3. Lernjahr) \\
\end{tabular}

\subsection{Voraussetzungen (Phase 0: VORWISSEN)}
Die SuS bringen aus dem 1. und 2. Lernjahr (Qué pasa 1) folgende Kenntnisse mit:
\begin{itemize}
    \item \textbf{Grammatik:} ser/estar (Unterscheidung), gustar/encantar (Dativ-Konstruktion), Präsens regelmäßiger und wichtiger unregelmäßiger Verben
    \item \textbf{Wortschatz:} Familie, Beschreibungen, Adjektive für Charakter und Aussehen
    \item \textbf{Kommunikativ:} Personen beschreiben, Vorlieben ausdrücken
\end{itemize}

\subsection{Differenzierung (intern)}

\textbf{WICHTIG:} Keine Niveaumarkierungen auf Schülermaterial!

\begin{tabular}{|l|p{3.5cm}|p{3.5cm}|p{3.5cm}|}
\hline
& \textcolor{basisgreen}{\textbf{[B] Basis}} & \textcolor{standardblue}{\textbf{[S] Standard}} & \textcolor{erweiterungpurple}{\textbf{[E] Erweiterung}} \\
\hline
\textbf{Ziel} & Berufsreife & Mittlere Reife & Gymnasium \\
\hline
\textbf{parecer} & Lückentext mit Vorgaben & Sätze mit Hilfe bilden & Freie Meinungsäußerung \\
\hline
\textbf{conocer} & Konjugation einsetzen & Sätze mit/ohne \textit{a} & Eigene Beispiele, Erklärung \\
\hline
\textbf{Hilfen} & Konjugationstabelle, Wortschatzliste & Satzanfänge & Nur Impulse \\
\hline
\end{tabular}

\subsection{Besonderheiten}
\begin{itemize}
    \item \textbf{LRS:} 2 SuS (Nachteilsausgleich: Zeitzugabe, größere Schrift auf AB)
    \item \textbf{DaZ:} 1 SuS (deutsche Erklärungen verstärken)
    \item \textbf{Leistungsstark:} 4 SuS (Zusatzaufgaben: eigene Dialoge)
\end{itemize}

% ============================================================================
% 2. SACHANALYSE
% ============================================================================
\section{Sachanalyse}

\subsection{Sprachliche Strukturen}

\subsubsection{Das Verb \textit{parecer} (scheinen, vorkommen)}
\begin{itemize}
    \item \textbf{Konstruktion:} Wie \textit{gustar} -- indirektes Objektpronomen + parecer + Adjektiv/Adverb
    \item \textbf{Konjugation:} me/te/le/nos/os/les + parece (Sg.) / parecen (Pl.)
    \item \textbf{Verwendung:} Meinungen über Personen, Dinge, Vorschläge äußern
    \item \textbf{Beispiele:}
    \begin{itemize}
        \item \textit{¿Qué te parece mi hermano?} -- \textit{Me parece simpático.}
        \item \textit{¿Qué os parecen los deberes?} -- \textit{Nos parecen difíciles.}
        \item \textit{¿Qué te parece quedar a las ocho?} -- \textit{Me parece bien.}
    \end{itemize}
\end{itemize}

\subsubsection{Das Verb \textit{conocer} (kennen, kennenlernen)}
\begin{itemize}
    \item \textbf{Besonderheit:} Unregelmäßige 1. Person Singular: \textbf{conozco} (c $\rightarrow$ zc vor o)
    \item \textbf{Konjugation:} conozco, conoces, conoce, conocemos, conocéis, conocen
    \item \textbf{conocer + a:} Bei Personen IMMER mit Präposition \textit{a}!
    \begin{itemize}
        \item \textit{¿Conoces \textbf{a} María?} (Person)
        \item \textit{¿Conoces la ciudad?} (Sache -- ohne \textit{a})
    \end{itemize}
    \item \textbf{Bedeutung:} ``kennen'' oder ``kennenlernen'' je nach Kontext
\end{itemize}

\subsection{Kontrastive Betrachtung (Deutsch $\leftrightarrow$ Spanisch)}
\begin{itemize}
    \item \textbf{parecer:} Im Deutschen keine direkte Entsprechung der Dativ-Konstruktion (``Mir scheint er nett'' ist veraltet $\rightarrow$ ``Ich finde ihn nett'')
    \item \textbf{conocer + a:} Im Deutschen keine Präposition bei Personen (``Ich kenne Maria'' vs. ``Conozco \textbf{a} María'') -- häufige Fehlerquelle!
    \item \textbf{Positiver Transfer:} Konstruktion ähnlich wie \textit{gustar} bereits bekannt
\end{itemize}

\subsection{Kulturelle Aspekte}
\begin{itemize}
    \item \textbf{Patchwork-Familien:} Text ``La nueva familia'' thematisiert moderne Familienstrukturen
    \item \textbf{Jugendsprache:} Ausdrücke wie \textit{¡Cómo mola!}, \textit{¡Qué fuerte!}, \textit{Es un palo} (S. 20)
    \item \textbf{Cádiz:} Handlungsort -- älteste Stadt Westeuropas, Andalusien
\end{itemize}

% ============================================================================
% 3. DIDAKTISCHE ANALYSE
% ============================================================================
\section{Didaktische Analyse}

\subsection{Stellung in der Sequenz}
Dies ist die \textbf{1. Doppelstunde} der Sequenz ``\reihenthema'' (4 Doppelstunden gesamt).

\begin{tabular}{|c|l|l|}
\hline
\textbf{DS} & \textbf{Thema} & \textbf{Schwerpunkt} \\
\hline
\textbf{01} & \textbf{parecer y conocer} & \textbf{Grammatik: parecer, conocer + a} \\
\hline
02 & La familia de Clara & Leseverstehen, Wortfeld Familie \\
\hline
03 & ser vs. estar (Vertiefung) & Grammatik: Zustand vs. Eigenschaft \\
\hline
04 & Mi familia -- Proyecto & Sprechfertigkeit, Präsentation \\
\hline
\end{tabular}

\subsection{Kompetenzziele (Rahmenplan MV)}
\begin{itemize}
    \item \textbf{Kommunikativ (Sprechen):} Meinungen über Personen äußern
    \item \textbf{Kommunikativ (Lesen):} Einen längeren Dialog verstehen (T2, S. 18)
    \item \textbf{Sprachlich (Grammatik):} parecer (Dativ-Konstruktion), conocer (yo-Form, + a)
    \item \textbf{Interkulturell:} Familienstrukturen in Spanien, Jugendsprache
\end{itemize}

\subsection{Legitimation}
\textbf{Rahmenplan Spanisch MV (Kl. 7-10):}
\begin{itemize}
    \item Kompetenzbereich: Sprechen -- an Gesprächen teilnehmen
    \item Standard: ``Meinungen, Gefühle und Wünsche äußern''
    \item Kompetenzbereich: Verfügen über sprachliche Mittel -- Grammatik
    \item Standard: ``Verben mit besonderen Konstruktionen anwenden''
\end{itemize}

% ============================================================================
% 4. STUNDENZIELE
% ============================================================================
\section{Stundenziele}

\subsection{Grobziel}
Die SuS erweitern ihre \textbf{kommunikative Kompetenz} im Bereich \textbf{Sprechen}, indem sie mit \textit{parecer} Meinungen über Personen äußern und mit \textit{conocer} über Bekanntschaften sprechen.

\subsection{Feinziele (operationalisiert)}

\begin{tabular}{|c|p{7cm}|c|c|}
\hline
\textbf{Nr.} & \textbf{Die SuS können...} & \textbf{AFB} & \textbf{Diff.} \\
\hline
\textbf{FZ1} & die Bedeutung von \textit{parecer} und \textit{conocer} im Kontext erschließen & I & alle \\
\hline
\textbf{FZ2} & \textit{parecer} mit indirektem Objektpronomen + Adjektiv verwenden, um Meinungen zu äußern & II & alle \\
\hline
\textbf{FZ3} & \textit{conocer} korrekt konjugieren (insb. \textit{conozco}) und die Regel \textit{conocer + a} bei Personen anwenden & II & alle \\
\hline
\textbf{FZ4} & in einem Partnerdialog Meinungen über Familienmitglieder austauschen & II & [S]/[E] \\
\hline
\textbf{FZ5} & die Unterscheidung Person (+ a) vs. Sache (ohne a) bei \textit{conocer} begründen & III & [E] \\
\hline
\end{tabular}

% ============================================================================
% 5. METHODISCHE ANALYSE
% ============================================================================
\section{Methodische Analyse}

\subsection{Phasierung (Doppelstunde = 2 × 45 Min.)}

\textbf{1. Stunde: Texterschließung + parecer}
\begin{enumerate}
    \item \textbf{Einstieg (5'):} Bildimpuls Patchwork-Familie -- Vorwissen aktivieren
    \item \textbf{Erarbeitung I (15'):} Text ``La nueva familia'' lesen (S. 18), globales Verstehen
    \item \textbf{Erarbeitung II (10'):} \textit{parecer} im Text markieren, Regel ableiten
    \item \textbf{Übung (10'):} Partnerarbeit -- Meinungen austauschen mit \textit{parecer}
    \item \textbf{Sicherung (5'):} Ergebnisse im Plenum, Merkkasten
\end{enumerate}

\textbf{2. Stunde: conocer + Anwendung}
\begin{enumerate}
    \item \textbf{Wiederholung (5'):} Schnelles Quiz zu \textit{parecer}
    \item \textbf{Erarbeitung III (10'):} \textit{conocer} im Text finden, yo-Form, Regel + a
    \item \textbf{Übung (15'):} Arbeitsblatt -- Lückentext conocer, conocer + a vs. ohne a
    \item \textbf{Anwendung (10'):} Partnerdialog -- ``¿Conoces a...?'' / ``¿Qué te parece...?''
    \item \textbf{Sicherung (5'):} Zusammenfassung, Hausaufgabe
\end{enumerate}

\subsection{Medien und Materialien}
\begin{itemize}
    \item Tafel/Whiteboard (Tafelanschrieb)
    \item Lehrwerk: \lehrwerk, \lehrwerkseite
    \item \textbf{AB-01-parecer-conocer.pdf} (Arbeitsblatt, 2 Seiten, 26 BE)
    \item \textbf{LP-01-parecer-conocer.html} (Lernpfad, optional digital)
    \item Bildimpuls: Familie (Beamer/Folie)
\end{itemize}

% ============================================================================
% 6. VERLAUFSPLANUNG
% ============================================================================
\section{Verlaufsplanung}

\textbf{1. Stunde (45 Min.) -- Texterschließung + parecer}

\begin{longtable}{|p{0.8cm}|p{1.3cm}|p{4.5cm}|p{3.5cm}|p{1cm}|p{1.5cm}|}
\hline
\textbf{Zeit} & \textbf{Phase} & \textbf{Lehrerhandeln} & \textbf{Schülerhandeln} & \textbf{SF} & \textbf{Medien} \\
\hline
\endhead

0--5 & Einstieg & Zeigt Bild einer Patchwork-Familie: ``¿Qué veis? ¿Cómo es una familia moderna?'' & Beschreiben, Vermutungen äußern & PL & Bild/ Beamer \\
\hline

5--20 & Erar\-bei\-tung I & Leitet Textarbeit an: ``Leed el texto T2 (S. 18). ¿De qué hablan los amigos?'' Aufgabe 10a (Themen zuordnen) & Lesen still, dann laut (Rollen), Themen identifizieren & EA, dann GA & Buch S. 18 \\
\hline

20--30 & Erar\-bei\-tung II & ``Buscad las formas de \textit{parecer}. ¿Cómo funciona?'' Sammelt an Tafel, leitet Regel ab & Markieren im Text, nennen Beispiele & PL & Tafel \\
\hline

30--40 & Übung & Partnerarbeit anleiten: ``¿Qué te parece Gonzalo / Clara / Nicolás?'' Dialogkarten verteilen & Dialoge üben, Meinungen austauschen & PA & Dialog\-karten \\
\hline

40--45 & Siche\-rung & Ergebnisse abfragen, Merkkasten diktieren & Vorstellen, notieren im Heft & PL & Tafel \\
\hline
\end{longtable}

\vspace{0.5cm}

\textbf{2. Stunde (45 Min.) -- conocer + Anwendung}

\begin{longtable}{|p{0.8cm}|p{1.3cm}|p{4.5cm}|p{3.5cm}|p{1cm}|p{1.5cm}|}
\hline
\textbf{Zeit} & \textbf{Phase} & \textbf{Lehrerhandeln} & \textbf{Schülerhandeln} & \textbf{SF} & \textbf{Medien} \\
\hline
\endhead

0--5 & Wieder\-holung & Schnelles Quiz: ``¿Qué te parece el español?'' -- SuS antworten & Antworten mit parecer & PL & -- \\
\hline

5--15 & Erar\-bei\-tung III & ``Buscad \textit{conocer} en el texto. ¿Qué es especial?'' Tafelanschrieb: conozco, conocer + a & Suchen, yo-Form erkennen, Regel + a verstehen & EA/ PL & Buch S. 18, Tafel \\
\hline

15--30 & Übung & Arbeitsblatt verteilen, Aufgaben 1--4 bearbeiten lassen, differenziert unterstützen & AB bearbeiten: Konjugation, + a einsetzen, eigene Sätze & EA & AB \\
\hline

30--40 & Anwen\-dung & Partnerdialog anleiten: ``Preguntad: ¿Conoces a...? ¿Qué te parece...?'' & Dialoge führen, beide Verben kombinieren & PA & -- \\
\hline

40--45 & Ab\-schluss & Zusammenfassung an Tafel, HA erklären & Notieren, Fragen stellen & PL & Tafel \\
\hline
\end{longtable}

\textbf{Legende:} EA = Einzelarbeit, PA = Partnerarbeit, GA = Gruppenarbeit, PL = Plenum

% ============================================================================
% 7. ERWARTETE SCHWIERIGKEITEN
% ============================================================================
\section{Erwartete Schwierigkeiten}

\begin{tabular}{|p{4.5cm}|p{8cm}|}
\hline
\textbf{Schwierigkeit} & \textbf{Reaktion} \\
\hline
\textit{parecer} vs. \textit{gustar} verwechselt & Parallele Gegenüberstellung an Tafel: ``Me gusta = ich mag / Me parece = ich finde'' \\
\hline
\textit{conocer + a} vergessen bei Personen & Merkregel: ``Personas con \textbf{a}!'' -- Visualisierung mit Person-Symbol \\
\hline
\textit{conozco} falsch als *\textit{conoco} & Explizit: c $\rightarrow$ zc vor o (wie andere Verben: traducir $\rightarrow$ traduzco) \\
\hline
Textverständnis T2 zu schwierig & Vorentlastung Wortschatz, Text in Abschnitten lesen \\
\hline
Heterogenität im Lerntempo & Differenzierte AB-Aufgaben: [B] mit Hilfen, [E] freie Produktion \\
\hline
Zeitproblem & Puffer: Partnerdialog kürzen, HA erweitern \\
\hline
\end{tabular}

% ============================================================================
% 8. TAFELANSCHRIEB
% ============================================================================
\section{Tafelanschrieb}

\begin{center}
\fbox{\parbox{14cm}{
\textbf{\large parecer y conocer}

\vspace{0.5cm}

\textbf{1. parecer} (wie \textit{gustar}!)

\begin{tabular}{ll}
me parece / me parecen & -- mir scheint / ich finde \\
te parece / te parecen & -- dir scheint / du findest \\
le parece / le parecen & -- ihm/ihr scheint \\
\end{tabular}

\textit{¿Qué te parece Gonzalo?} -- \textit{Me parece simpático.}

\vspace{0.5cm}

\textbf{2. conocer} (kennen / kennenlernen)

\begin{tabular}{llllll}
conozco & conoces & conoce & conocemos & conocéis & conocen \\
\end{tabular}

\vspace{0.3cm}

\textbf{conocer + a} bei \textbf{Personen}!

\textit{¿Conoces \textbf{a} María?} -- Sí, conozco \textbf{a} María.

\textit{¿Conoces la ciudad?} -- Sí, conozco la ciudad. (Sache = ohne \textit{a})
}}
\end{center}

% ============================================================================
% 9. HAUSAUFGABE
% ============================================================================
\section{Hausaufgabe}

\begin{itemize}
    \item \textbf{Pflicht:} Lehrwerk S. 21, Aufgabe 14b (conocer-Lücken ergänzen)
    \item \textbf{Vokabeln lernen:} parecer, conocer, engreído, simpático, tranquilo, cansado
    \item \textbf{Optional [E]:} 5 Sätze über eigene Familie mit \textit{parecer} und \textit{conocer}
    \item \textbf{Optional:} Lernpfad LP-01-parecer-conocer.html bearbeiten
\end{itemize}

% ============================================================================
% 10. MATERIALIEN / ANLAGEN
% ============================================================================
\section{Anlagen}

\begin{enumerate}
    \item \textbf{AB-01-parecer-conocer.pdf} -- Arbeitsblatt (2 Seiten, 26 BE)
    \item \textbf{ML-01-parecer-conocer.pdf} -- Musterlösung (Lösungen in rot)
    \item \textbf{LH-01-parecer-conocer.pdf} -- Lehrerhinweise + Kurzplan
    \item \textbf{LP-01-parecer-conocer.html} -- Lernpfad (digital, optional)
    \item Dialogkarten für Partnerarbeit (Kopiervorlage in LH)
    \item Bildimpuls: Patchwork-Familie
\end{enumerate}

% ============================================================================
% 11. REFLEXIONSFRAGEN
% ============================================================================
\section{Reflexionsfragen (für Lehrkraft)}

\begin{enumerate}
    \item Haben alle SuS die Regel \textit{conocer + a} verstanden?
    \item War die Textarbeit zu T2 zeitlich angemessen?
    \item Wurden die Differenzierungsangebote genutzt?
    \item Konnten die SuS im Partnerdialog beide Verben anwenden?
    \item Welche Fehler traten gehäuft auf? $\rightarrow$ Folgestunde berücksichtigen
\end{enumerate}

% ============================================================================
% 12. DIDAKTIK-SCORE
% ============================================================================
\newpage
\section{Didaktik-Score}

\begin{scorebox}
\textbf{Bewertung der didaktischen Qualität nach 10 Dimensionen}\\
\textbf{Ziel-Score: $\geq$ 9.0 (Premium-Qualität)}

\vspace{1em}
\begin{tabular}{|c|l|c|c|p{4.5cm}|}
\hline
\textbf{\#} & \textbf{Dimension} & \textbf{Gew.} & \textbf{Score} & \textbf{Notiz} \\
\hline
1 & Differenzierung & 15\% & 9/10 & [B]/[S]/[E] qualitativ verschieden \\
\hline
2 & Taxonomie (AFB) & 10\% & 9/10 & AFB I: 20\%, AFB II: 60\%, AFB III: 20\% \\
\hline
3 & Lernziel-Alignment & 10\% & 10/10 & RP MV Sprechen + Grammatik \\
\hline
4 & Aktivierung & 10\% & 9/10 & PA-Dialoge, aktive Textarbeit \\
\hline
5 & Scaffolding & 10\% & 9/10 & Gestufte Hilfen, Merkkasten \\
\hline
6 & Feedback-Qualität & 10\% & 8/10 & Im LP elaboriert, AB Selbstkontrolle \\
\hline
7 & Sprachliche Klarheit & 10\% & 10/10 & Altersgerecht, klare Anweisungen \\
\hline
8 & Motivation \& Relevanz & 10\% & 9/10 & Alltagsbezug Familie, Jugendsprache \\
\hline
9 & Inklusion (UDL) & 10\% & 9/10 & Visuell (Tafel), auditiv (Dialog), schriftlich \\
\hline
10 & Praktikabilität & 5\% & 9/10 & 90 Min. realistisch, Puffer vorhanden \\
\hline
\multicolumn{3}{|r|}{\textbf{GESAMT:}} & \textbf{9.1/10} & \textcolor{successgreen}{\textbf{FREIGABE}} \\
\hline
\end{tabular}

\vspace{1em}
\textbf{Bewertungsschwellen:}
\begin{itemize}
    \item[$\boxtimes$] \textcolor{successgreen}{$\geq$ 9.0: \textbf{FREIGABE} (Premium-Qualität)}
    \item[$\Box$] \textcolor{warningorange}{8.0--8.9: Iteration empfohlen}
    \item[$\Box$] 6.0--7.9: Überarbeitung erforderlich
    \item[$\Box$] $<$ 6.0: Grundlegende Neukonzeption
\end{itemize}
\end{scorebox}

\end{document}
